\section{Methods}
\label{sec:methods}

\subsection{Design and scientific data flow}

\ROBERT{} separates six stages: observation ingestion, atmosphere construction,
opacity evaluation, radiative transfer, instrument response, and inference.
Typed immutable containers carry pressure and spectral grids, stellar and
planetary properties, atmospheric state, observations, and spectra. A retrieval
parameter vector is therefore distinct from the physical atmosphere that it
generates. This separation is intended to make conventions---including units,
grid orientation, reference radius, opacity identity, and interpolation
policy---visible at component boundaries.

For data set $d$, the forward path is
\begin{equation}
 \boldsymbol{\theta}\rightarrow
 \mathcal{A}(\boldsymbol{\theta})\rightarrow
 \mathcal{K}\rightarrow
 \mathcal{R}\rightarrow
 \mathcal{I}_d\rightarrow
 \boldsymbol{m}_d(\boldsymbol{\theta}),
 \label{eq:dataflow}
\end{equation}
where $\mathcal{A}$ constructs the atmosphere, $\mathcal{K}$ evaluates opacity,
$\mathcal{R}$ solves radiative transfer, and $\mathcal{I}_d$ applies the data
set's response and binning. Opacity preparation is cached outside likelihood
calls. Run manifests record the model configuration, priors, opacity and input
checksums, software environment, random seed, and grid hashes.

\subsection{Atmospheric structure and composition}

The atmosphere is represented on an explicitly ordered pressure grid. Available
temperature prescriptions include isothermal, tabulated, Guillot-type
\citep{Guillot2010}, Madhusudhan--Seager-type \citep{MadhusudhanSeager2009},
and Parmentier--Guillot-type profiles \citep{ParmentierGuillot2014}.
Composition may be supplied as free constant-with-altitude mixing ratios,
layer-dependent profiles, or equilibrium chemistry through FastChem
\citep{Stock2018}. Trace-gas abundances are combined with an explicit H$_2$/He
background, with validation of non-negative abundances and normalization.

Under hydrostatic equilibrium, the absorber column in layer $i$ is computed as
\begin{equation}
 N_i = \frac{\Delta P_i}{\bar{\mu}_i m_{\rm u} g_i},
 \label{eq:column}
\end{equation}
where $\Delta P_i$ is the pressure difference, $\bar{\mu}_i$ the mean molecular
weight in atomic mass units, $m_{\rm u}$ the atomic mass constant, and $g_i$
the gravitational acceleration. Spherical shell radii can be obtained by
integrating hydrostatic balance from a declared reference pressure and radius.

\subsection{Opacity and correlated-$k$ treatment}

\ROBERT{} reads NEMESIS KTA tables and opacity archives prepared with
\code{exo\_k} \citep{Leconte2021}. Gas opacity is interpolated in pressure and
temperature under an explicit bounds policy and evaluated on the model's
spectral and $g$ grids. For gas $s$, layer $i$, wavelength bin $j$, and ordinate
$l$, the monochromatic-distribution optical depth is
\begin{equation}
 \tau_{sijl}=k_{sijl}\,x_{si}\,N_i,
 \label{eq:tau}
\end{equation}
with unit conversion determined by opacity metadata. Multiple gases may share
an ordinate grid or be combined by random overlap with resorting and
rebinning. Collision-induced absorption and Rayleigh extinction are separate
contributors, allowing their provenance and approximations to be inspected.

For random overlap, two discrete opacity distributions with optical depths
$\{\tau_a\}$ and $\{\tau_b\}$ and quadrature weights $\{w_a\}$ and $\{w_b\}$
form the product distribution
\begin{equation}
 \tau_{ab}=\tau_a+\tau_b, \qquad w_{ab}=w_aw_b.
 \label{eq:random-overlap}
\end{equation}
The $n_g^2$ product points are ranked and conservatively rebinned to the
original $n_g$ intervals; the procedure is applied successively to all active
absorbers. The reference implementation explicitly sorts the product
distribution. In the accelerated implementation, nondecreasing input
distributions---the normal correlated-$k$ case---are recognized and the
$n_g$ already sorted rows of outer sums are combined with an $n_g$-element
heap. Rebinning is performed in the same streaming pass with reusable work
arrays. This avoids a general sort of every $n_g^2$ temporary at every layer
and wavelength. Unsorted inputs retain the general sorting path, and a gas is
skipped only when its maximum optical depth is below the declared transparent
gas threshold. The NumPy implementation remains the scientific reference;
the optional compiled backend is required to reproduce it to floating-point
precision.

\subsection{Profile-guided computational optimization}
\label{sec:performance-methods}

Optimization used warmed forward-model profiles of an 80-layer, six-molecule
WASP-69b calculation with 16 $g$ ordinates, evaluated on both a 1591-point
native opacity grid and 280-point observation grids. Changes were accepted
only if they preserved the layer, wavelength, angular, and correlated-$k$
resolution, equations, convergence thresholds, and double-precision
arithmetic. Consequently, the changes described here alter implementation
cost rather than retrieval physics.

Random overlap was the dominant initial cost and was optimized as described
above. Two additional repeated operations were removed from likelihood calls.
First, prepared opacity objects retain validated spectral lookup indices, and
providers may cache $\log k$ for immutable coefficient tables instead of
recomputing logarithms during every pressure--temperature interpolation. The
cache is explicit and optional because it exchanges memory for speed; cached
and uncached calculations are tested for exact equality. Second, four fixed
four-term source contractions in the SH4 solver were written as explicit sums,
avoiding thousands of dispatches to a general tensor-contraction routine while
leaving the spherical-harmonics equations unchanged. Spherical transmission
was not rewritten because profiling showed that its contribution was already
small.

The random-overlap benchmark spans one, three, and six species at 280 and 1591
wavelengths, records throughput and memory, and compares a deterministic
subset with the NumPy reference on every run. On the development system
(macOS arm64, Python 3.12.13, NumPy 2.2.6, and Numba 0.61.2), the largest
observed random-overlap difference was $3.13\times10^{-13}$ in optical depth.
The complete 279-test suite passes after these changes. In the representative
WASP-69b workloads, the native-grid-then-bin path decreased from 6.35 to
1.21\,s per warmed evaluation (a factor of 5.3), while preparing opacities per
instrument mode decreased from 1.34 to 0.43\,s (a factor of 3.1). The
corresponding log-likelihoods were unchanged exactly within the recorded
double-precision calculation. Cached log-coefficients used approximately
106\,MiB for the four observation-grid providers and 599\,MiB for the full
native-grid provider, so production configurations should select the cache
policy with available memory in mind.

\subsection{Thermal emission and transmission}

For a non-scattering, plane-parallel ray with direction cosine $\mu$, the
top-of-atmosphere intensity is evaluated from the formal solution
\begin{equation}
 I_{\lambda}(0,\mu) = I_{\lambda}(\tau_{\rm b},\mu)
 e^{-\tau_{\rm b}/\mu} +
 \int_{0}^{\tau_{\rm b}} B_{\lambda}[T(t)]
 e^{-t/\mu}\frac{\mathrm{d}t}{\mu}.
 \label{eq:formal}
\end{equation}
When temperatures are specified at layer edges, \ROBERT{} integrates a Planck
source that is linear in optical depth within each layer; otherwise it records
the constant layer-centre approximation. Disc integration supports normal
emission, a uniform thermal disc, and Lobatto quadrature. Eclipse depths divide
the planet flux by either an input stellar spectrum or a blackbody stellar
reference.

Clouds carry extinction optical depth $\Delta\tau$, single-scattering albedo
$\omega_0$, and asymmetry factor $g$. A grey mass-extinction coefficient
$\kappa_{\rm cld}$ gives
\begin{equation}
 \Delta\tau_{{\rm cld},i}=0.1\,
 \kappa_{\rm cld}\frac{\Delta P_i}{g_i},
 \label{eq:grey-cloud}
\end{equation}
where $\kappa_{\rm cld}$ is in cm$^2$ g$^{-1}$ and the factor $0.1$ converts to
SI units. Multiple scattering can be solved with a hemispheric-mean two-stream
method based on \citet{Toon1989} or a four-term spherical-harmonics (SH4)
method following \citet{Rooney2023}, including Henyey--Greenstein moments and
optional delta-$M$ scaling. The current transmission path integrates slant
optical depths through hydrostatic spherical shells; \roberttodo{state the
validated refraction, scattering, and cloud limitations for the release used
in the paper}.

\subsection{Inhomogeneous daysides and multiple data sets}

The emergent spectrum of a two-region dayside is mixed after radiative transfer,
\begin{equation}
 F_{\lambda}=f_{\rm hot}F_{\lambda,\rm hot}+
 (1-f_{\rm hot})F_{\lambda,\rm cold},
 \label{eq:two-region}
\end{equation}
where the columns may have independent thermal and cloud structures while
sharing selected chemistry parameters. This ordering is essential because
radiative transfer and the Planck function are nonlinear. A diluted one-region
model is the limiting case in which the complementary region contributes
negligible flux. Each observation retains its instrument identity, wavelength
bins, mask, and nuisance parameters; the joint likelihood is the sum of the
per-data-set terms.

\subsection{Likelihoods and inference}

For independent Gaussian uncertainties, the log likelihood is
\begin{equation}
 \ln \mathcal{L}=-\frac{1}{2}\sum_d\sum_i
 \left[\frac{(y_{di}-m_{di})^2}{s_{di}^2}
 +\ln(2\pi s_{di}^2)\right],
 \quad s_{di}^2=\sigma_{di}^2+\sigma_{\rm jit,d}^2.
 \label{eq:likelihood}
\end{equation}
The same sampler-independent retrieval problem can be evaluated with optimal
estimation \citep{Rodgers2000} or nested sampling. The current Bayesian adapter
uses UltraNest's reactive nested sampler \citep{Buchner2021}; priors and their
unit-cube transforms are explicit objects. Model comparisons will report
evidence differences only after convergence and repeated-run checks.

\subsection{Numerical verification and benchmarking}
\label{sec:benchmarks}

Validation is layered. Unit tests cover grids, conservation and limiting cases,
opacity interpolation, optical-depth assembly, radiative-transfer limits,
likelihoods, manifests, and inference adapters. Injection--recovery tests then
exercise the complete parameter-to-spectrum path. Cross-code comparisons use
matched atmospheric states and, where possible, identical optical-depth arrays
to isolate the radiative-transfer kernel from opacity differences.

For a controlled 160-layer comparison, matched \ROBERT{} and PICASO SH4 point
radiances agree to a maximum relative difference of $3.23\times10^{-6}$ for
both isotropic scattering and a case with $\omega_0=0.9$ and $g=0.6$. Using
real molecular correlated-$k$ structure for HAT-P-32b, the maximum and RMS
disc-integrated differences are currently $0.0507$ and $0.0151$ per cent,
respectively; a FastChem plus collision-induced-absorption case reaches a
maximum difference of $0.00361$ per cent. \roberttodo{Regenerate these values
from a tagged release and add a table containing grid, hardware, package
versions, metric definitions, and output-file hashes.}

The petitRADTRANS benchmark suite compares stable reference calculations for
multispecies emission and transmission and tests convergence with native HDF5
opacities. \roberttodo{Insert the final petitRADTRANS version, exact common
inputs, residual metrics, and figure after freezing the reference data.}

\subsection{WASP-69b end-to-end benchmark}

The WASP-69b benchmark uses the 280-point published $2.455$--$11.875\,\mu$m
spectrum of \citet{Schlawin2024}, retaining NIRCam/F322W2, the NIRCam overlap,
NIRCam/F444W, and MIRI/LRS as four named data sets. We compare one-region clear,
two-region clear, diluted one-region, and two-region grey-cloud models. The
primary validation quantities are the recovered hot-region fraction, thermal
profiles, cloud contrast, metallicity, C/O, maximum likelihood, and evidence
relative to the published analysis. \roberttodo{Run and report the staged
retrieval suite; the current repository contains the data model and target
workflow but not final publication-grade posteriors.}
